

\begin{lemma}
\label{lem:braid_invariance_property_2}
Let $(W, S)$ be a Coxeter system. Consider a group element $w \in W$ and two words of finite length related by a braid move  $\word{i} = \word{i}_1 (ij \ldots) \word{i}_2 \to \word{i}_1 (j i \ldots) \word{i}_2 = \word{i}'$.
Then $\word{i}$ does not contain any subwords which are non-reduced expressions of 

$w$
 if and only if $\word{i}'$ does not contain any subwords which are non-reduced expressions of 
 
 $w$.
 \end{lemma}



\begin{proof}
 If we have a non-reduced  subword for $w$ in $\word{i}'$ which is not a subword of $\word{i}$, there are two cases. Either it contains the entire consecutive subword $(ji \ldots)$, in which case the same positions in $\word{i}$ give a non-reduced  subword for $w$. 
 Or it contains a piece of $(ji \ldots)$ with some letters removed in the middle, which  has either some $ii$ with the letter $j$ between them omitted, or some $j j$ with the letter $i$ between them omitted. 
 If there are multiple such pairs, delete one, because it will still be a non-reduced subword for $w$. 
 This may create a new such pair, but the length of the subword will be shorter. 
 Repeat until there is exactly one such pair and after deleting it no new pair is created.  
 Except for this pair, the part of the remaining subword inside $(ji \ldots)$ has length at most $m_{ij} - 3$ and is either an initial or a final subword of $(i j \ldots)$. 
 Attach to it either the last and third to last letters of $(ij \ldots)$, if it was initial, or first and third letters of $(ij \ldots)$, if it was final.
We get a non-reduced subword for $w$ inside $\word{i}$. 
Thus, if the property did not hold for $\word{i}'$, it does not hold for $\word{i}$ either. The other implication follows by switching the roles of $\word{i}$ and $\word{i}'$.
\end{proof}

\begin{proposition}
\label{prop:no_nonreduced_double_root_free}
Let $(W, S)$ be a Coxeter system, $\word{i}$ a word of finite length. Then $\word{i}$ does not contain any subwords which are non-reduced expressions of $\dem := \dem(\word{i})$ if and and only if $\word{i}$ is double root free.
\end{proposition}

\begin{proof}
 If the subword $\word{j}$ has the same Demazure product as $\word{i}$, then each facet of the subword complex of $\word{j}$ extends to a facet of the subword complex of $\word{i}$ by adding letters in the complement $\word{i} \backslash \word{j}$. The root configuration for the former is a subset of the root configuration for the latter. Thus, $\word{i}$ is double root free if and only if each subword with the same Demazure product is double root free. The other condition (about non-reduced subwords), by 
 its definition, also holds for $\word{i}$ if and only if it holds for each subword of $\word{i}$ with the same Demazure product $\dem$.

 We now show that we can reduce the proof to the case of words of length $\ell(\dem) +2$. First, if $\word{i}$ is not double root free, the spherical subword complex $\Delta(\word{i}, \dem(\word{i}))$ has a facet $I$ having two (automatically flippable) positions $a, b$ with the same root in the root configuarion of $I$.  Then the subword $\word{j}$ of $\word{i}$ formed by $\word{i} \backslash I, a, b$ is not double root free.  It has length $\ell(\dem) + 2$, and $\dem(\word{j}) = \dem$. 
 Second, if $\word{i}$ does contain a non-reduced subword of $\dem$, then by the deletion property \cite[Proposition 1.4.7]{bb_coxeter}, we can remove pairs of letters of this subword until we get a subword of length $\ell(\dem) + 2$ still giving a non-reduced expression of $\dem$. Thus, each of the two conditions holds for $\word{i}$ if and only if it holds for each subword of $\word{i}$ of length $\ell(\dem) + 2$ with Demazure product $\dem$, and so it is sufficient to prove the proposition for words of length $\ell(\dem) + 2$. 



Assume now that we have a word $\word{j}$ of length $\ell(\dem) + 2$ and with $\dem(\word{j}) = \dem$. 

Then $\word{j}$ has the form $\word{i}(u) i \word{i}(v) j \word{i}(w)$, with $\word{i}(u)\word{i}(v)\word{i}(w)$ being a reduced expression of $\dem$. Then the root configuration for the face of the subword complex which is a complement for this reduced expression is $\{u(\alpha_i), uv(\alpha_j)\}$. We have $\pi(\word{i}(u)i \word{i}(v) j \word{i}(w)) = \dem = \pi(\word{i}(u)\word{i}(v)\word{i}(w))$ if and only if $\pi(i \word{i}(v) j \word{i}(w)) = \pi(\word{i}(v)\word{i}(w))$, if and only if 
$\pi(i \word{i}(v) j) = \pi(\word{i}(v))$, if and only if $\pi(i \word{i}(v)) = \pi(\word{i}(v) j)$, if and only if $s_i (v (\alpha_j)) = v(s_j(\alpha_j))$, if and only if $s_i(v(\alpha_j)) = - v(\alpha_j)$, if and only if $v(\alpha_j) = \pm \alpha_i$, if and only if $uv(\alpha_j) = \pm u\alpha_i$. 


Now that we proved that the failure of  first condition is equivalent to the word being ``not double root free, up to sign'', the result would follow if we prove that this implies the failure of the honest double root free condition.
 Assume that we have $\pi\left(\word{i}(u) i \word{i}(v) j \word{i}(w)\right) = uvw$, equivalently,  $v(\alpha_j) = \pm \alpha_i$. Since roots for the letters in the complement to our facet are in bijection with the inversions of $uvw$, 
 there is at most one letter in this complement which has the same root up to sign. Since the positions with letters $i, j$ are flippable, such letter must exist.  

This letter then can be replaced by $i$ and by $j$ by (two different) flips, and since no other flips of the maximal face $\{i, j\}$ are possible while the complex is spherical, the subword complex is a triangle formed by this letter $k$ and the positions of 
$i$ and $j$, i.e. the maximal simplices are $\{i, j\}, \{i, k\}, \{k, j\}$. 
By \cite[Lemma 2.6]{BergeronCeballos} (extending \cite[Lemmas 3.3 and 3.6]{cls} to possibly infinite Coxeter groups),
both roots in the facet formed by the two leftmost of these positions are positive, so the root configuration for  this leftmost (called \emph{positive greedy} in \cite{PilaudStump-EL}) facet 
consists of two copies of the same positive root. This completes the proof.
\end{proof}

\begin{corollary} 
\label{cor:drf_braid_invariant}
Let $(W, S)$ be a Coxeter system. Consider two words
of finite length related by a braid move $\word{i} = \word{i}_1 (ij \ldots) \word{i}_2 \to \word{i}_1 (j i \ldots) \word{i}_2 = \word{i}'$. Then $\word{i}$ is double root free if and only if $\word{i}'$ is double root free.
\end{corollary}